In this section, we describe the sources and rationale underlying the parameter values and prior distributions used in the model. The complete list of parameter definitions, numerical values, calibration priors, and references is provided in Table~\ref{tab-params}, with the sensitivities of the screening approaches reported separately in Table~\ref{tab-screening-sens}. Parameters were informed using a combination of published literature, previous modelling studies, empirical observations from Kiribati, and calibration to epidemiological data. Where direct evidence was unavailable or substantial uncertainty existed, prior distributions were specified to reflect plausible ranges and updated through Bayesian calibration.

\subsection{TB natural history parameters}

The model incorporates a detailed representation of the spectrum of \textit{M. tuberculosis} infection and disease. Parameters governing progression through early infection states were informed by our previous work on TB latency parameterisation, which estimated parameters for different model structures using TB progression data [REFERENCE PLACEHOLDER]. While the model structure used in this analysis differs from those previously calibrated, we can establish a correspondence between key parameters based on their role within the early infection process.

The previous parameterisation study calibrated several alternative latent infection structures. The comparison presented here focuses on Model 4, which represented early infection using two latent compartments. Individuals entered an initial latent state and could either progress directly to active TB or transition to a secondary latent state associated with slower progression dynamics. The parameter \(\epsilon\) governed direct progression from the initial latent state to disease, while the parameter \(\kappa\) governed transition from the initial latent state into the secondary latent state.

In the present model, newly infected individuals enter an incipient infection state (\(E\)), from which they may either progress to active disease or transition to a contained infection state (\(C\)). Contained infection may subsequently undergo immune-mediated clearance to \(X\) or return to the incipient state following immunological breakdown.

The progression parameter in the present model, denoted \(\epsilon_i\), governs direct progression from the initially occupied infection state (\(E\)) to active disease. This parameter therefore has the same interpretation as \(\epsilon\) in Model 4, namely the rate at which newly infected individuals progress directly to disease before entering a more controlled infection state. This correspondence can be demonstrated analytically by considering disease incidence immediately following infection. When all individuals are initialised in the high-risk infection state, the incidence rate at \(t = 0\) reduces to the direct progression parameter in both model structures, establishing the equivalence of \(\epsilon\) and \(\epsilon_i\) with respect to early progression dynamics. Consequently, estimates of \(\epsilon\) from the previous calibration were used directly to inform age-specific values of \(\epsilon_i\).

No direct mathematical correspondence exists between the containment rate in the present model (\(\gamma_i\)) and the transition rate \(\kappa\) estimated in Model 4. The present formulation includes additional pathways, notably immune breakdown from the contained state back to the incipient state, which were not represented in the earlier model. Nevertheless, \(\kappa\) provides the most relevant source of information because both parameters govern movement from an initially infected state with elevated risk of disease progression into a state associated with greater immune control and lower short-term progression risk. Estimates of \(\kappa\) were therefore used to inform \(\gamma_i\) in the present model.

Importantly, the revised model allows the balance between rapid progression and longer-term latent infection to be determined jointly by the containment rate (\(\gamma_i\)) and the immune breakdown rate (\(\omega\)). Calibration of the breakdown parameter therefore provides flexibility to reproduce a range of early infection trajectories that could not be represented within the simpler Model 4 structure.

The remaining parameter governing early infection dynamics, the clearance rate (\(\xi\)), describes spontaneous resolution of contained infection and has no direct analogue in the previous model. This parameter was assigned a prior distribution reflecting the plausible range of spontaneous clearance rates reported in the literature [REFERENCE PLACEHOLDER] and updated through calibration.

\textcolor{red}{NEED TO CHECK/REVISE FROM HERE ONWARDS}

Similarly, parameters governing susceptibility to reinfection following containment, infection clearance, or recovery from disease were assigned prior distributions reflecting current uncertainty in the literature and were updated through calibration.

The active disease component distinguishes individuals according to both clinical status (subclinical versus clinical disease) and infectiousness (less infectious versus more infectious disease). Parameters governing transitions between these disease states were selected to reproduce realistic distributions of disease phenotypes while remaining identifiable from the available data. To reduce identifiability issues, transition rates from more infectious to less infectious disease and from clinical to subclinical disease were fixed at values corresponding to an average duration of approximately one year in each state before transition. Transition rates in the opposite direction were estimated during calibration.

Relative infectiousness parameters were informed by the emerging literature on subclinical TB and heterogeneity in transmission potential [REFERENCE PLACEHOLDER]. Individuals classified as less infectious were assumed to generate 40\% of the transmission produced by more infectious individuals, while subclinical disease was assumed to be half as infectious as clinical disease. The combined infectiousness of each disease category was determined by multiplying the corresponding relative infectiousness factors.

Parameters governing TB mortality, self-recovery, and treatment relapse were informed by published estimates and previous modelling studies [REFERENCE PLACEHOLDER]. Where substantial uncertainty remained, prior distributions were assigned and updated through calibration.

The model additionally includes a transmission-scaling parameter that allows a continuum between frequency-dependent and density-dependent transmission. This parameter was introduced to account for uncertainty regarding the relationship between population density and transmission intensity in Kiribati, where increasing population density and limited habitable land may plausibly influence contact patterns and transmission risk. The parameter was assigned a uniform prior spanning the full range between purely frequency-dependent and purely density-dependent transmission.

\subsection{Programmatic parameters}

The structural representation of passive detection, treatment and screening is described in Section~\ref{sec:screening}, and that of screening reachability in Section~\ref{sec:reachability}. This section documents the values assigned to the associated parameters, the evidence used to inform them, and the reasons for estimating some during calibration while fixing others.

\subsubsection{Passive detection}

No direct estimate of the passive detection rate was available for Kiribati. The recent detection level (\(\delta_{\mathrm{rec}}\)), the inflection time of the scale-up (\(t^\ast\)) and the past-to-recent ratio (\(\phi\)) were therefore assigned uninformative uniform priors and estimated during calibration against the historical notification series (Table~\ref{tab-params}). The prior on \(\delta_{\mathrm{rec}}\) spans two orders of magnitude, corresponding to an average delay from onset of clinical disease to treatment initiation ranging from around two months to around ten years, and is therefore not expected to constrain the posterior.

The steepness parameter \(\sigma\) was fixed at 0.15 per year rather than estimated. The timing and the abruptness of the scale-up are only weakly separable given a single notification time series, and allowing both to vary introduced strong correlation between the two parameters without improving the fit.

The relative detection rate for subclinical disease was fixed at zero and the relative detection rate in the hard-to-reach stratum at 0.5, following the assumptions set out in Sections~\ref{sec:screening} and~\ref{sec:reachability}. Neither quantity is identifiable from the available data.

\subsubsection{Treatment}

The average treatment duration \(\tau\) was fixed at 0.5 years, corresponding to the standard six-month first-line regimen.

The treatment success rate \(\mathrm{TSR}(t)\) was taken directly from the annual treatment outcomes reported for Kiribati [REFERENCE PLACEHOLDER] and implemented as a linear interpolation between the reported values, which have ranged between 83\% and 95\% since 1995. Outcomes were not reported before that date, so the series was extended backwards from an assumed value of zero in 1945, preceding the availability of effective chemotherapy, to 50\% in 1950. These historical values are not intended as estimates of programme performance but serve to prevent the model from attributing an implausible treatment benefit to the pre-programmatic era; the calibration targets are insensitive to them.

The fraction of negative outcomes resulting in death, \(f_{\mathrm{neg}}\), was fixed at 40\%, approximating the proportion of unsuccessful outcomes recorded as deaths in the treatment outcomes reported to WHO for Kiribati [REFERENCE PLACEHOLDER]. The remaining unsuccessful outcomes, comprising treatment failure, loss to follow-up and cases not evaluated, are returned to active disease in the model. This proportion fluctuates markedly between years because the annual cohorts are small, so a single representative value was used rather than a time-varying series, and it was not estimated during calibration.

\subsubsection{Reachability parameters}

The screening-reachable fraction \(\pi\) was fixed at 0.85, based on participation observed during the completed phase of the PEARL study. The relative susceptibility of the hard-to-reach stratum (\(\theta_{\mathrm{sus}} = 1.5\)) and its relative passive detection rate (\(\theta_{\mathrm{det}} = 0.5\)) are not informed by the available data and could not be estimated jointly with the transmission and detection parameters, as they act on the same flows. They were therefore fixed at plausible values, and \(\theta_{\mathrm{sus}}\) was varied in sensitivity analysis.

\subsubsection{Screening algorithm sensitivities}

The sensitivity assigned to each screening algorithm for each model state is reported in Table~\ref{tab-screening-sens}. With the exception of the TST parameters, these values were derived from the completed phase of the PEARL campaign, in which each participant found to have TB was classified by clinical status and by infectiousness, and the result of each nested algorithm was recorded. Sensitivity was evaluated against all disease detected by the full symptom-screening, CXR and Xpert algorithm among participants with an interpretable Xpert result. That algorithm therefore has a sensitivity of one for every disease state by construction, because it defines the reference standard, and this is the value assigned to the Xpert-containing branch in the model. It should be noted that this branch is not the same as the delivered PEARL programme, which combines it with CXR alone in the proportions given in Section~\ref{sec:screening}.

The sensitivity of CXR to subclinical disease was the only screening quantity that was materially uncertain, and it was estimated during calibration. Uniform priors were specified as the exact (Clopper--Pearson) 95\% confidence intervals around the sensitivities observed in each subclinical stratum: 4 of 16 individuals with less-infectious subclinical disease were identified (25\%, 95\% CI 7--52\%), against 44 of 52 with more-infectious subclinical disease (85\%, 95\% CI 72--93\%). The considerably wider prior for the less-infectious state reflects the small number of individuals in that stratum rather than any additional biological uncertainty. These two parameters are further informed during calibration by the CXR-measured prevalence target, which is compared with a model output constructed using the same sensitivities, alongside the prevalence measured under the full algorithm.

The probability of a positive TST was fixed at 0.75 for both the incipient and contained infection states [REFERENCE PLACEHOLDER]. The corresponding probability among individuals who have cleared infection or recovered from disease was estimated during calibration, since the persistence of tuberculin reactivity after clearance is poorly characterised; this parameter is informed by the TST positivity measured during the PEARL screening phase and, in combination with the infection-state values, determines the mapping between modelled infection prevalence and observed TST positivity.

The probability of completing preventive treatment once identified as TST-positive was fixed at 0.7, consistent with completion levels observed during implementation [REFERENCE PLACEHOLDER]. Because completion is a key determinant of the benefit attributable to the TPT arm of the campaign, a sensitivity analysis was run assuming a completion probability of 0.6.

\subsubsection{Screening coverage}

Coverage levels of 65\%, 75\% and 85\% of the total population were examined. These are scenario assumptions rather than estimated quantities: the lower two represent plausible degrees of incomplete participation, while the highest corresponds to screening 99.4\% of the screening-reachable stratum and is treated as the practical upper bound achievable without extending the campaign to the hard-to-reach population (Section~\ref{sec:reachability}).
